% 图契约：本题的源、边界和推导逐项对应；箭头符号按图注解释。
\documentclass[tikz,border=7pt]{standalone}
\usepackage{fontspec,xeCJK,amsmath,amssymb}
\IfFontExistsTF{Times New Roman}{\setmainfont{Times New Roman}}{\setmainfont{TeX Gyre Termes}}
\IfFontExistsTF{Heiti SC}{\setCJKmainfont[AutoFakeBold=1.5]{Heiti SC}}{\setCJKmainfont[AutoFakeBold=1.5]{FandolHei-Regular}}
\usetikzlibrary{arrows.meta,bending,calc,patterns,positioning}
\definecolor{ink}{HTML}{243745}
\definecolor{blue}{HTML}{275DAB}
\definecolor{green}{HTML}{177C63}
\definecolor{red}{HTML}{BD4935}
\definecolor{orange}{HTML}{AA6C16}
\definecolor{purple}{HTML}{715393}
\tikzset{every node/.style={font=\fontsize{9.3}{12.5}\selectfont,text=ink},
axis/.style={-{Stealth[length=4pt]},line width=.65pt,draw=ink!65},
vec/.style={-{Stealth[length=5pt]},line width=1pt,draw=blue},
short/.style={-{Stealth[length=3pt]},line width=.8pt},
guide/.style={dashed,draw=ink!40,line width=.5pt},
card/.style={rounded corners=2pt,draw=blue!65,fill=blue!4,align=center,text width=6.5cm,inner xsep=9pt,inner ysep=7pt},
title/.style={font=\fontsize{10.5}{14}\selectfont\bfseries,text width=15cm,align=center},
label/.style={font=\fontsize{8.4}{11}\selectfont,align=center},
note/.style={label,text width=6.2cm}}
\begin{document}
\begin{tikzpicture}[x=1cm,y=1cm]
\path[use as bounding box] (-8,-3.5) rectangle(8,3.5);
\node[title] at (0,3.15) {M3-033\quad 点电荷与介质平面的力和相互作用能};
\begin{scope}[shift={(-4,0)}]
\fill[blue!6] (-2.4,-1.8) rectangle(2.4,0);\draw[ink,thick] (-2.4,0)--(2.4,0);
\node[label] at(-1.25,2.05){a区：真空};\node[label] at(-1.65,-1.65){b区：$\varepsilon_r$};
\draw[axis] (-2.15,.2)--(-2.15,1.25) node[left] {$z$};
\draw[<->,orange] (-.75,0)--(-.75,1.15) node[midway,left] {$h$};
\draw[<->,orange] (-.75,-1.15)--(-.75,0) node[midway,left] {$h$};
\node[label] at(1.7,.2){$z=0$};
\fill[red] (0,1.15) circle(2.3pt) node[above=3pt] {$Q$};
\draw[blue,dashed] (0,-1.15)circle(3pt);\node[right] at(.1,-1.15){$Q'$（像源）};
\draw[vec,red] (0,.95)--(0,.3) node[right] {$\mathbf F$};
\node[note] at(0,-2.5){图示取 $Q>0,\ \varepsilon_r>1$；力指向介质。\\$Q'$ 只在 a 区用作反射像源。};\end{scope}
\node[card,draw=blue!65,fill=blue!4] at (3.9,1.8) {$\displaystyle\mathbf F=\frac{QQ^\prime}{16\pi\varepsilon_0h^2}\mathbf e_z$};
\node[card,draw=green!65,fill=green!4] at (3.9,0) {$\displaystyle U_{\rm int}=\frac12Q\varphi_{\rm ind}=\frac{QQ^\prime}{16\pi\varepsilon_0h}$};
\node[card,draw=orange!65,fill=orange!4] at (3.9,-1.8) {$\displaystyle -\frac{dU_{\rm int}}{dh}=F_z$};
\end{tikzpicture}\end{document}